\documentclass[12pt,a4paper]{report}

% -------------------------
% Packages
% -------------------------
\usepackage[utf8]{inputenc}
\usepackage[T1]{fontenc}
\usepackage[french]{babel}

\usepackage{geometry}
\geometry{
    left=3cm,
    right=2.5cm,
    top=2.5cm,
    bottom=2.5cm
}

\usepackage{setspace}
\onehalfspacing

\usepackage{graphicx}
\usepackage{float}
\usepackage{caption}
\usepackage{subcaption}

\usepackage{amsmath, amssymb}
\usepackage{booktabs}
\usepackage{array}
\usepackage{longtable}

\usepackage{xcolor}
\usepackage{hyperref}
\usepackage{url}

\usepackage{fancyhdr}
\usepackage{titlesec}
\usepackage{enumitem}

% -------------------------
% Configuration
% -------------------------
\hypersetup{
    colorlinks=true,
    linkcolor=black,
    urlcolor=blue,
    citecolor=black,
    % pdftitle={Titre du rapport},
    % pdfauthor={Nom Prénom}
}

\pagestyle{fancy}
\fancyhf{}
\fancyhead[L]{Neurosymbolisme}
\fancyhead[R]{Pier-Luc Dupéré}
\fancyfoot[C]{\thepage}

\setlength{\parindent}{0pt}
\setlength{\parskip}{6pt}


% -------------------------
% Début du document
% -------------------------
\begin{document}

% -------------------------
% Page de garde
% -------------------------
\begin{titlepage}

    \begin{center}

        \includegraphics[width=0.50\textwidth]{img/logo-universite.png}
        
        \vspace{1cm}

        {\Large \textbf{Université du Québec à Rimouski}}\\[0.3cm]
        {\large Département de mathématiques, d'informatique et de génie}\\[0.3cm]
        {\large 8INF976 - Sujet spécial en intelligence artificielle}\\[1.5cm]

        \rule{\textwidth}{0.5pt}\\[0.4cm]

        {\Huge \textbf{Neurosymbolisme}}\\[0.4cm]
        {\Large Rapport de lecture}\\[0.4cm]

        \rule{\textwidth}{0.5pt}\\[2cm]

        \begin{flushleft}
            \textbf{Réalisé par :} Pier-Luc Dupéré \\[0.4cm]
            \textbf{Professeur :} Mehdi Adda, PhD 
        \end{flushleft}

        \vfill
    \end{center}

\end{titlepage}

\tableofcontents
% \listoffigures
% \listoftables

% -------------------------
% Corps du rapport
% -------------------------
\clearpage
\pagenumbering{arabic}


% -------------------------
% Chapitre 1
% -------------------------
\chapter{Unifying Large Language Models and Knowledge Graphs: A Roadmap[1]}

\section{Mise en contexte}
Malgré leur succès dans de nombreux domaines, les grands modèles de langage (LLM) 
sont parfois critiqués pour leur manque de fiabilité. Même s'ils apprennent des 
faits à partir des données utilisées pour leur entraînement. Ils peuvent les oublier, 
les déformer ou inventer des informations. Ce phénomène est appelé « hallucination ».

les grands modèles de langage (LLM) ont montré de grandes capacités dans de nombreuses 
tâches de traitement automatique du langage naturel la plupart d'entre eux utilisent 
l'architecture Transformer, composée de :
\begin{itemize}
    \item Langage à encodeur seul
    \item Langage à encodeur-décodeur
    \item Langage à décodeur seul
\end{itemize}

Afin de remédier à ce problème, les chercheurs ont proposé d'intégrer des graphes de 
connaissances (KG) dans les LLM. Ces graphes permettent d'organisent les informations sous forme de triplets.
Ce triplet est composé de l'entité de tête, la relation et l'entité de queue.
(head entity, relation, tail entity)

Les graphes de connaissances existants peuvent être classés en quatre groupes selon le type d'informations 
qu'ils contiennent :
\begin{itemize}
    \item Les KG encyclopédiques
    \item Les KG de sens commun
    \item Les KG spécifiques à un domaine
    \item Les KG multimodaux
\end{itemize}

Les graphes de connaissances permettent un raisonnement symbolique et produisent des résultats 
plus faciles à comprendre. L'association des LLM et des graphes de connaissances intéresse de 
plus en plus les chercheurs, car ces deux technologies peuvent se compléter et s'améliorer 
mutuellement.

Dans cet article, trois méthodes sont proposées pour intégrer les graphes de connaissances 
dans les LLM : l'enrichissement des LLM par les graphes de connaissances, l'amélioration des 
graphes de connaissances grâce aux LLM.

% ----------------------------------------------
\section{Les LLM enrichis par des KG}
Les grands modèles de langage (LLM) sont généralement entraînés de manière non supervisée sur 
de vastes corpus. Malgré leurs bonnes performance. Ils peuvent manquer de connaissances 
pratiques et actualisées sur le monde réel.

Les graphes de connaissances (KG) peuvent être intégrés aux LLM de trois principales façons :
\begin{itemize}
    \item En les intégrant à l'objectif d'entraînement
    \item En ajoutant les connaissances des KG aux entrées du modèle
    \item En utilisant les KG lors de l'instruction-tuning
\end{itemize}

Il existe une problématique leur mise à jour nécessite d'un nouvel entraînement du modèle. 
Ces approches risquent donc de mal gérer des connaissances nouvelles ou absentes des données 
d'entraînement.

Les LLM sont souvent critiqués pour leur  difficulté à comprendre leur fonctionnement interne 
et leurs décisions. Les KG représentent les connaissances de manière structurée et explicite, 
ce qui peut rendre les raisonnements du modèle plus faciles à expliquer.

% ----------------------------------------------
\section{Les KG augmentés par des LLM}
(Embeddings) de KG augmentés par LLM : 
L'objectif est de convertir les entités et relations en vecteurs mathématiques. Contrairement aux méthodes classiques limitées à la structure du graphe.


Complétion de KG augmentée par LLM : 
Cette méthode vise à déduire les faits manquants dans un graphe. 
On utilise des pipelines de LLM pour prédire les descriptions d'entités et identifier l'élément manquant dans un triplet.

Construction de KG augmentée par LLM : 
Les LLM font la création de graphes structurés en automatisant la découverte d'entités, la résolution de coréférences et l'extraction de relations.

Génération de texte à partir de KG : 
Ce processus transforme des données structurées en texte naturel pour pallier le manque de 
données d'entraînement. On exploite les connaissances intrinsèques des LLM.

Questions-Réponses sur KG : 
Le but est de répondre à des questions naturelles via des faits structurés. 
Les LLM comblent le fossé entre le langage souple et le graphe rigide en jouant deux rôles.

\section{Les KG augmentés par des LLM}

\begin{itemize}
    \item \textbf{Embeddings de KG augmentés par LLM :} L'objectif est de convertir les entités et relations en vecteurs mathématiques.
    \item \textbf{Complétion de KG augmentée par LLM :} On utilise des pipelines de LLM pour prédire les descriptions d'entités et identifier l'élément manquant dans un triplet.
    \item \textbf{Construction de KG augmentée par LLM :} Les LLM automatisent la création de graphes via la découverte d'entités, la résolution de coréférences et l'extraction de relations.
    \item \textbf{Génération de texte à partir de KG :} Ce processus transforme des données structurées en texte naturel en exploitant les connaissances intrinsèques des LLM.
    \item \textbf{Questions-Réponses sur KG :} Les LLM permettent de répondre à des questions naturelles entre le langage utilisateur et les faits structurés du graphe.
\end{itemize}


% ----------------------------------------------
\section{La synergie entre LLM et KG}
Cette section présente la synergie entre les grands modèles de langage et les graphes de 
connaissances. Les LLM permettent de comprendre le langage naturel, tandis que les graphes 
fournissent des connaissances structurées et factuelles. Leur combinaison peut donc améliorer 
la représentation des connaissances et les capacités de raisonnement.

La synergie repose sur deux aspects principaux 
\begin{itemize}
    \item La représentation conjointe des connaissances 
    \item Le raisonnement combiné
\end{itemize}

\textbf{Représentation conjointe des connaissances} :
Les textes contiennent des informations souvent implicites et désorganisées, 
alors que les graphes de connaissances présentent des informations explicites 
et structurées. Plusieurs modèles combinent donc les deux sources grâce à des 
encodeurs et à des modules de fusion. Certains filtrent les informations inutiles 
du graphe, pendant que d'autres apprennent simultanément les représentations textuelles 
et les relations entre les entités. Cette approche permet de mieux exploiter les 
connaissances issues des textes et des graphes dans différentes tâches.

\textbf{Raisonnement combiné} :
Cette approche cherche à utiliser à la fois les informations textuelles et 
les relations présentes dans les graphes pour produire des réponses. 
Certains modèles utilisent deux encodeurs distincts qui échangent des 
informations, tandis que d'autres établissent une interaction plus fine 
entre les mots du texte et les entités du graphe. Des méthodes permettent 
également de réduire dynamiquement le graphe afin de se concentrer sur les 
informations les plus pertinentes.

Enfin, les deux approches présentent des avantages et des limites. La fusion des LLM 
et des graphes permet un entraînement conjoint et de bonnes performances. Cependant, 
elle augmente le nombre de paramètres et le coût de calcul. L'utilisation des LLM comme 
agents est plus flexible et ne nécessite pas toujours un entraînement supplémentaire. 
La définition de leurs actions et de leurs stratégies de recherche reste complexe. 
La combinaison des LLM et des graphes de connaissances demeure donc un domaine de 
recherche encore en développement.

% ----------------------------------------------
\section{Critiques et limites}
C'est une excellente synthèse critique. L'article contient plus de 250 citations qui sont pour 
la plus par ressente. De plus, il est intéressant de constater que deux approches sont possibles 
pour combiner les LLM et les graphes de connaissances (KG), afin d'obtenir des résultats plus 
précis et plus fiables. Ces deux approches sont complémentaires et peuvent permettre de compenser 
leurs limites respectives.

Cet article ouvre également plusieurs perspectives de recherche. Les auteurs présentent les 
orientations futures ainsi que plusieurs étapes importantes dans le développement de la 
recherche visant à unifier les graphes de connaissances et les LLM.
\begin{itemize}
    \item Détection des hallucinations
    \item Mise à jour des connaissances
    \item Injection de connaissances dans les LLM propriétaires
    \item Prise en compte de plusieurs modalités
    \item Compréhension de la structure des graphes
    \item Raisonnement bidirectionnel
\end{itemize}


% -------------------------
% Chapitre 2
% -------------------------
\chapter{From Local to Global: A GraphRAG Approach to Query-Focused Summarization [2]}

\section{Mise en contexte}
Le RAG est une méthode qui permet à un modèle de langage de rechercher des informations dans des sources 
externes avant de générer une réponse. Cette approche est utile lorsque les données sont trop nombreuses 
pour être intégrées entièrement dans la fenêtre de contexte du modèle.

Le RAG vectoriel récupère les documents dont le contenu est sémantiquement proche de la question grâce 
aux représentations vectorielles. Il est efficace pour retrouver des informations précises. Il est limité 
lorsqu'il faut comprendre les tendances et les relations générales présentes dans l'ensemble d'un corpus.

La GraphRAG améliore cette approche en construisant un graphe à partir des données. Elle regroupe les 
éléments liés en communautés thématiques.Par la suite, il génère des résumés à plusieurs niveaux qui vont faire 
des informations locales à une compréhension globale du corpus. Elle est donc particulièrement adaptée 
aux questions nécessitant une analyse d'ensemble. Ces graphes améliorent la recherche et fournissent 
une base solide aux réponses générées.

\section{Méthodologie de l'approche}

\textbf{Documents sources => Segments de texte}
Les documents du corpus sont d'abord découpés en plusieurs segments de texte. Le LLM analyse ensuite 
chaque segment afin d'en extraire les informations. La détermination de la taille des segments est essentielle.
Les coûts de traitement sont réduits avec des segments plus longs. Cependant, ils peuvent également diminuer 
la capacité du modèle à retrouver les informations dans segment.

\textbf{Segments de texte => Entités et relations}
Durant cette étape, le LLM est invité à extraire des instances d'entités importantes ainsi que les relations
entre ces entités à partir d'un segment de texte donné. De plus, le LLM génère de brèves descriptions
pour ces entités et relations.

\textbf{Entités et relations => Graphe des connaissances}
L'extraction des entités, des relations et des affirmations peut générer plusieurs occurrences d'un même 
élément dans différents documents. Lors de la construction du graphe de connaissances, ces occurrences 
sont regroupées pour créer des nœuds et des arêtes uniques. Les descriptions sont ensuite combinées et 
résumées. Fait important, le nombre d'occurrences d'une même relation détermine le poids de l'arête. 
Les affirmations sont regroupées de la même manière.

\textbf{Graphe de connaissances => Communautés du graphe}
Divers algorithmes de détection de communautés peuvent être utilisés pour partitionner le graphe en 
groupes de nœuds fortement connectés de manière hiérarchique. Cette méthode consiste à détecter récursivement 
des sous-communautés au sein de chaque communauté identifiée.


\textbf{Communautés du graphe => Résumés de communauté}
Cette étape consiste à générer un résumé pour chaque communauté de la hiérarchie. Les résumés permettent de 
mieux comprendre la structure globale et le contenu du corpus. Ils peuvent également être utilisés pour 
explorer les données sans requête précise. C'est en GraphRAG que ces résumés de communauté qui sont composés
d'éléments (noeuds, arêtes et affirmations associées)

\textbf{Résumés par communauté → Réponses communautaires → Réponse globale}
Les résumés des communautés sont d'abord répartis en plusieurs groupes. Le LLM produit ensuite une réponse 
pour chaque groupe et lui attribue un score d'utilité. Les réponses les plus pertinentes sont conservées. 
Puis, elles sont combinées et optimisées pour générer une réponse globale à la question de l'utilisateur.



\section{Analyses et discussions}
Les résultats montrent que les méthodes globales, en particulier GraphRAG, obtiennent de meilleurs résultats 
que le RAG vectoriel. GraphRAG est donc plus adapté aux questions nécessitant une compréhension globale d'un corpus.

GraphRAG présente aussi un avantage en matière de consommation de ressources. Les résumés des communautés 
racines nécessitent beaucoup moins de jetons que le traitement direct de tous les textes. Ce qui surprend 
c'est que les GraphRAG peut réduire considérablement la quantité de données envoyées au modèle tout en 
conservant de bonnes performances.

Les travaux futurs pourraient améliorer GraphRAG en combinant la recherche vectorielle avec la génération 
dynamique de résumés communautaires. Le système pourrait aussi fonctionner de manière hiérarchique. C'est à dire, 
qui'il pourrait passer d'une vue globale à une analyse plus détaillée des communautés.

Personnellement, je considère comme des points forts de cette article la présence d'un lien vers le dépôt Git, 
du pseudocode décrivant l'algorithme et des prompts détaillés présentés en annexe. Cela permet de reproduire les 
expériences et de mieux comprendre le fonctionnement de l'approche GraphRAG.

% ----------------------------------------------
\section*{Notes}
La rédaction et la mise en forme de ce document ont été réalisées avec \LaTeX{} et 
Visual Studio Code. Les outils de correction orthographique (Code Spell checker) et d'autocomplétion 
(IntelliSense) ont été utilisés comme aide à la rédaction.


% -------------------------
% Bibliographie
% -------------------------
\begin{thebibliography}{9}

\bibitem{site}
Pan, S., Luo, L., Wang, Y., Chen, C., Wang, J. et Wu, X.
\textit{Unifying Large Language Models and Knowledge Graphs: A Roadmap},
\url{https://arxiv.org/pdf/2306.08302},
consulté le 18 septembre 2026.

\bibitem{site}
Edge, D., Trinh, H., Cheng, N., Bradley, J., Chao, A., Mody, S., Truitt, S. et Larson, J. 
\textit{From Local to Global: A GraphRAG Approach to Query-Focused Summarization},
\url{https://arxiv.org/pdf/2404.16130v2},
consulté le 18 septembre 2026.

\end{thebibliography}

\end{document}
